\documentclass[11pt,letterpaper]{article}

\usepackage[T1]{fontenc}
\usepackage{lmodern}
\usepackage{microtype}
\usepackage[margin=1in]{geometry}
\usepackage{amsmath,amssymb,amsthm}
\usepackage[hidelinks]{hyperref}

\emergencystretch=2em

\newtheorem{theorem}{Theorem}[section]
\newtheorem{lemma}[theorem]{Lemma}

\newcommand{\N}{\mathbb{N}}
\newcommand{\Nzero}{\mathbb{N}_0}
\newcommand{\ZM}{\mathbb{Z}/M\mathbb{Z}}
\newcommand{\cpl}[1]{#1^{\mathsf{c}}}
\newcommand{\ints}[2]{[#1,#2]_{\mathbb{Z}}}

\title{Infinitely Many Aperiodic Greedy Pair-Sum-Avoiding Sequences}
\author{Zhiheng Li}
\date{}

\hypersetup{
  pdftitle={Infinitely Many Aperiodic Greedy Pair-Sum-Avoiding Sequences},
  pdfauthor={Zhiheng Li}
}

\begin{document}

\maketitle

\begin{abstract}
We construct an explicit infinite family of finite seeds whose greedy pair-sum-avoiding extensions have consecutive-gap sequences that are not eventually periodic. For every integer \(p\geq18\), the construction embeds the same scale-eight sumset controller in the residue classes \(p,2p,4p\) modulo \(7p\). A uniform finite shield gives a seed of cardinality \(7p-20\) and an infinite set \(S_p\) satisfying the exact greedy recurrence \(t\in S_p\Longleftrightarrow t\notin S_p+S_p\) beyond the seed.
\end{abstract}

\section{Introduction and the main theorem}

\noindent\textit{AI-use disclosure.} This proof was developed with assistance from OpenAI's GPT-5.6 Sol. All arguments were checked and verified by the author.

The question appears in Erd\H{o}s and Graham~\cite[p.~53]{ErdosGraham1980} and is catalogued as Erd\H{o}s Problem~341 in Bloom's online database~\cite{Bloom341}. Related periodicity criteria are due to Calkin and Finch~\cite{CalkinFinch1996}, while Calkin and Erd\H{o}s~\cite{CalkinErdos1996} proved that a natural class of aperiodic sum-free sets is incomplete. We answer the arbitrary finite-seed question negatively by constructing infinitely many counterexamples.

Write \(\N=\{1,2,3,\ldots\}\) and \(\Nzero=\{0,1,2,\ldots\}\). For \(X,Y\subseteq\Nzero\) and \(c\in\Nzero\), put
\[
X+Y=\{x+y\mid x\in X,\ y\in Y\},
\qquad
c+X=\{c+x\mid x\in X\},
\qquad
\cpl{X}=\Nzero\setminus X.
\]
Equal summands are allowed in every sumset.

Let \(A=\{a_1<\cdots<a_k\}\subseteq\N\) be a nonempty finite set. Install all elements of \(A\) as the initial terms. For \(n\geq k\), define
\begin{equation}\label{eq:greedy-definition}
a_{n+1}
=\min\{t\in\N\mid t>a_n,\ t\notin A_n+A_n\},
\qquad
A_n=\{a_1,\ldots,a_n\}.
\end{equation}
Let \(d_m=a_{m+1}-a_m\) denote the consecutive gaps.

Fix an integer \(p\geq18\), put \(M=7p\), and work in \(\ZM\), represented by \(\{0,\ldots,M-1\}\). Define
\begin{equation}\label{eq:parametric-PBRFD}
\begin{gathered}
P=\{p,2p,4p\},
\qquad
B=\{M-3,M-1,6p+3\},
\qquad R=P\cup B,\\
F=R\cup(R-B)\cup(R-P),
\qquad
D=\ZM\setminus F,
\end{gathered}
\end{equation}
where \(R-B=\{r-b:r\in R,\ b\in B\}\), and similarly for \(R-P\). Set
\begin{equation}\label{eq:seed-family}
A_p=D\cup B\cup(M+B)\cup\{p,4p,M+2p,M+4p\}.
\end{equation}

\begin{theorem}\label{thm:main}
For every integer \(p\geq18\), the seed \(A_p\) in~\eqref{eq:seed-family} has \(7p-20\) elements. For the sequence generated from \(A_p\) by~\eqref{eq:greedy-definition},
\[
\forall K\geq1\ \forall h\geq1\ \exists m\geq K
\quad d_{m+h}\neq d_m.
\]
Thus its consecutive-gap sequence is not eventually periodic. In particular, the sets \(A_p\), \(p\geq18\), form an infinite family of counterexamples.
\end{theorem}

\begin{lemma}[Exact recurrence implies the greedy rule]\label{lem:greedy}
Let \(S\subseteq\N\) be infinite, let \(N\in S\), and put \(A=S\cap\ints{1}{N}\). Suppose that
\begin{equation}\label{eq:exact-recurrence}
t\in S
\quad\Longleftrightarrow\quad
t\notin S+S
\qquad (t>N).
\end{equation}
Then the increasing enumeration of \(S\) is exactly the greedy extension of \(A\).
\end{lemma}

\begin{proof}
Scan the integers \(t>N\) in order. Inductively, immediately before testing \(t\), the selected set is \(S\cap\ints{1}{t-1}\). Since all elements of \(S\) are positive, every representation \(t=u+v\) with \(u,v\in S\) has \(u,v<t\), and hence
\[
t\in(S\cap\ints{1}{t-1})+(S\cap\ints{1}{t-1})
\quad\Longleftrightarrow\quad t\in S+S.
\]
Thus~\eqref{eq:exact-recurrence} makes the rule select exactly the elements of \(S\), completing the induction.
\end{proof}

\section{A scale-eight sumset controller}

For integers \(a\leq b\), write \(\ints{a}{b}=\{a,a+1,\ldots,b\}\). We use the standard identity
\begin{equation}\label{eq:interval-sum}
\ints{a}{b}+\ints{c}{d}=\ints{a+c}{b+d}
\qquad (a\leq b,\ c\leq d).
\end{equation}

Define
\begin{align}
L&=\bigcup_{k\geq0}\ints{8^k}{4\cdot8^k-1}, \\
Y_0&=\{0\}\cup
      \bigcup_{k\geq0}\ints{4\cdot8^k}{8\cdot8^k-1},
      \label{eq:Y0}\\
Y_1&=\ints{1}{3}\cup
      \bigcup_{k\geq0}\ints{16\cdot8^k-1}{32\cdot8^k-1},
      \label{eq:Y1}\\
Y_2&=\ints{0}{1}\cup\ints{7}{15}\cup
      \bigcup_{k\geq0}\ints{64\cdot8^k-1}{128\cdot8^k-1}.
      \label{eq:Y2}
\end{align}
At each scale, the corresponding blocks of \(L\) and \(Y_0\setminus\{0\}\) partition \(\ints{8^k}{8^{k+1}-1}\). Thus
\begin{equation}\label{eq:Y0-complement-L}
Y_0=\cpl{L}.
\end{equation}

\begin{lemma}[Controller identities]\label{lem:controller}
The sets in~\eqref{eq:Y0}--\eqref{eq:Y2} satisfy
\begin{align}
Y_1&=\cpl{(Y_0+Y_0)},                                      \label{eq:gate01}\\
Y_2&=\cpl{(Y_1+Y_1)},                                      \label{eq:gate12}\\
Y_0&=\{0\}\cup\bigl(1+\cpl{(Y_2+Y_2)}\bigr).             \label{eq:gate20}
\end{align}
\end{lemma}

\begin{proof}
A direct interval calculation gives
\begin{align}
Y_0+Y_0
&=\{0\}\cup\bigcup_{k\geq0}
  \ints{4\cdot8^k}{16\cdot8^k-2}=\cpl{Y_1},
  \label{eq:Y0-square}\\
Y_1+Y_1
&=\ints{2}{6}\cup\bigcup_{k\geq0}
  \ints{16\cdot8^k}{64\cdot8^k-2}=\cpl{Y_2},
  \label{eq:Y1-square}\\
Y_2+Y_2
&=\ints{0}{2}\cup\ints{7}{30}\cup\bigcup_{k\geq0}
  \ints{64\cdot8^k-1}{256\cdot8^k-2}=L-1.
  \label{eq:Y2-square}
\end{align}
Here \(L-1=\{\ell-1\mid\ell\in L\}\).
Equations~\eqref{eq:Y2-square} and~\eqref{eq:Y0-complement-L} yield~\eqref{eq:gate20}.
\end{proof}

\begin{lemma}\label{lem:Y0-aperiodic}
The membership indicator of \(Y_0\) is not ultimately periodic.
\end{lemma}

\begin{proof}
Suppose that there are \(N\geq0\) and \(h\geq1\) such that
\[
n\in Y_0\iff n+h\in Y_0
\qquad(n\geq N).
\]
Choose \(a=8^k>\max\{N,h\}\) and put \(m=4a-h\). Then \(N<a\leq m\leq4a-1\). Hence \(m\in L\), so \(m\notin Y_0\), whereas \(m+h=4a\in\ints{4a}{8a-1}\subseteq Y_0\). This contradicts the assumed periodicity.
\end{proof}

\section{A uniform finite shield modulo \texorpdfstring{\(7p\)}{7p}}

All residues in this section lie in \(\ZM\), with the notation of~\eqref{eq:parametric-PBRFD}.

\begin{lemma}[Uniform finite shield]\label{lem:shield}
In $\ZM$, for every \(p\geq18\),
\begin{align}
D+B&=\ZM\setminus R,                                      \label{eq:cover-mod}\\
R\cap\bigl((B+B)\cup(P+B)\cup(D+B)\cup(D+P)\bigr)
&=\varnothing,                                             \label{eq:avoid-mod}\\
(P+P)\cap R&=P.                                            \label{eq:PP-mod}
\end{align}
Moreover, the only ordered pairs in \(P\times P\) whose sum lies in \(R\) are
\begin{equation}\label{eq:diagonal-pairs}
p+p=2p,
\qquad
2p+2p=4p,
\qquad
4p+4p=M+p.
\end{equation}
\end{lemma}

\begin{proof}
Expanding the definition in~\eqref{eq:parametric-PBRFD} gives
\begin{equation}\label{eq:F-explicit}
\begin{split}
F=\{&0,2,\ p-6,p-4,p,p+1,p+3,\ 2p-3,2p,2p+1,2p+3,\\
&3p-3,3p-1,3p,\ 4p,4p+1,4p+3,\\
&5p-3,5p-1,5p,5p+3,\\
&6p-3,6p-1,6p,6p+3,6p+4,6p+6,\ M-3,M-2,M-1\}.
\end{split}
\end{equation}
For \(p\geq18\), these are \(30\) distinct representatives in \(\{0,\ldots,M-1\}\). In particular, \(|D|=M-30\). The definition of \(F\) immediately gives
\begin{equation}\label{eq:direct-avoidance}
(D+B)\cap R=\varnothing,
\qquad
(D+P)\cap R=\varnothing.
\end{equation}

We prove the reverse inclusion in~\eqref{eq:cover-mod}. Let \(r\in\ZM\), and put \(f=r+1\). If \(f\notin F\), then \(f\in D\) and
\[
r=f+(M-1)\in D+B.
\]
It remains to consider \(f\in F\). Put
\[
G=\{0,p+1,2p+1,4p+1,6p+4,M-2\},
\qquad
E=\{p-6,3p-3,5p-3,6p-3,M-3\}.
\]
If \(f\in G\), then direct subtraction shows that \(r=f-1\in R\). If \(f\in F\setminus(G\cup E)\), comparison with~\eqref{eq:F-explicit} shows that \(f+2\notin F\), and hence
\[
r=(f+2)+(M-3)\in D+B.
\]
For the five values of \(f\) in \(E\), the corresponding representations are
\begin{align*}
p-7&=(2p-10)+(6p+3),&
3p-4&=(4p-7)+(6p+3),\\
5p-4&=(6p-7)+(6p+3),&
6p-4&=(M-7)+(6p+3),\\
M-4&=(p-7)+(6p+3).
\end{align*}
Thus every residue outside \(R\) belongs to \(D+B\), and~\eqref{eq:direct-avoidance} proves~\eqref{eq:cover-mod}.

Direct calculation shows that \(B+B\) and \(P+B\) are disjoint from \(R\). Together with~\eqref{eq:direct-avoidance}, this proves~\eqref{eq:avoid-mod}. Finally, the three mixed sums of elements of \(P\) are \(3p,5p,6p\), all outside \(R\), while the diagonal sums are those in~\eqref{eq:diagonal-pairs}. This proves~\eqref{eq:PP-mod} and the ordered-pair assertion.
\end{proof}

We now define the infinite set
\begin{equation}\label{eq:S}
\begin{split}
S_p={}&D\cup\{Mq+b\mid q\in\Nzero,\ b\in B\}\\
&\cup\{Mq+p\mid q\in Y_0\}
 \cup\{Mq+2p\mid q\in Y_1\}
 \cup\{Mq+4p\mid q\in Y_2\}.
\end{split}
\end{equation}
The initial controller memberships give
\begin{equation}\label{eq:seed-prefix}
S_p\cap\ints{1}{2M-1}
=D\cup B\cup(M+B)\cup\{p,4p,M+2p,M+4p\}=A_p.
\end{equation}
In particular, \(2M-1=M+(M-1)\in A_p\). The four sets in~\eqref{eq:seed-prefix} are disjoint, and therefore
\begin{equation}\label{eq:seed-size}
|A_p|=(M-30)+3+3+4=M-20=7p-20.
\end{equation}

\begin{theorem}\label{thm:recurrence}
For every \(p\geq18\) and every \(t>2M-1\),
\begin{equation}\label{eq:main-recurrence}
t\in S_p
\quad\Longleftrightarrow\quad
t\notin S_p+S_p.
\end{equation}
\end{theorem}

\begin{proof}
Write \(t=Mq+r\), with \(0\leq r<M\). Since \(t>2M-1\), we have \(q\geq2\).

First let \(r\in R=B\cup P\). The six possible pair-sum types are
\[
D+D,\quad D+B,\quad D+P,\quad B+B,\quad B+P,\quad P+P.
\]
A \(D+D\) sum is at most \(2M-2<t\). Equation~\eqref{eq:avoid-mod} excludes the next four types from \(R\), while~\eqref{eq:PP-mod} and~\eqref{eq:diagonal-pairs} leave only the three diagonal \(P+P\) pairs. Hence
\begin{align*}
Mq+2p\in S_p+S_p&\iff q\in Y_0+Y_0,\\
Mq+4p\in S_p+S_p&\iff q\in Y_1+Y_1,\\
Mq+p\in S_p+S_p&\iff q-1\in Y_2+Y_2.
\end{align*}
Indeed, the corresponding representations are
\begin{align*}
(Mu+p)+(Mv+p)&=M(u+v)+2p,\\
(Mu+2p)+(Mv+2p)&=M(u+v)+4p,\\
(Mu+4p)+(Mv+4p)&=M(u+v+1)+p.
\end{align*}
The exclusion of every other source type proves the converse implication in each equivalence. The controller identities~\eqref{eq:gate01}--\eqref{eq:gate20} therefore make every \(P\)-rail integer represented exactly when it is absent from \(S_p\). Every \(B\)-rail integer belongs to \(S_p\), and the same exhaustive source analysis shows that none is represented. Thus~\eqref{eq:main-recurrence} holds for \(r\in R\).

Now let \(r\notin R\). Such an integer \(t\) is not in \(S_p\): its residue excludes every infinite rail, and \(t>2M-1\) excludes the finite set \(D\). By~\eqref{eq:cover-mod}, there exist \(d\in D\) and \(b\in B\) such that
\[
d+b=Mc+r
\qquad(c\in\{0,1\}).
\]
Since \(q\geq2\), we have \(q-c\geq1\), and
\[
t=d+\bigl(M(q-c)+b\bigr)\in S_p+S_p.
\]
The second summand is a positive element of a full \(B\)-rail. Both summands are positive and sum to \(t\), so each is smaller than \(t\). This completes the proof.
\end{proof}

By~\eqref{eq:seed-prefix}, Theorem~\ref{thm:recurrence}, and Lemma~\ref{lem:greedy}, the increasing enumeration of \(S_p\) is exactly the greedy extension of \(A_p\).

\section{Nonperiodicity of the gaps}

\begin{lemma}\label{lem:gaps-imply-membership}
Let \(T=\{t_1<t_2<\cdots\}\subseteq\N\) be infinite. If the consecutive gaps \(t_{m+1}-t_m\) are eventually periodic, then the membership indicator of \(T\) is ultimately periodic.
\end{lemma}

\begin{proof}
Suppose that the gaps have period \(h\geq1\) from index \(K\) onward. Let
\[
H=\sum_{j=0}^{h-1}(t_{K+j+1}-t_{K+j})>0.
\]
Every block of \(h\) consecutive gaps in the periodic tail has sum \(H\), so \(t_{m+h}=t_m+H\) for \(m\geq K\). Since \(T\) is infinite, it is unbounded. Thus, given \(n\geq t_K\), there is a unique \(m\geq K\) such that \(t_m\leq n<t_{m+1}\), and
\[
t_{m+h}\leq n+H<t_{m+h+1},
\qquad
(n+H)-t_{m+h}=n-t_m.
\]
Thus \(n=t_m\) if and only if \(n+H=t_{m+h}\), which proves
\[
n\in T\iff n+H\in T
\qquad(n\geq t_K).
\]
This is ultimate periodicity of the membership indicator.
\end{proof}

\begin{proof}[Proof of Theorem~\ref{thm:main}]
Equation~\eqref{eq:seed-size} proves the asserted cardinality. The residue-\(p\) rail of \(S_p\) satisfies
\begin{equation}\label{eq:rail-p}
Mq+p\in S_p\quad\Longleftrightarrow\quad q\in Y_0
\qquad(q\in\Nzero).
\end{equation}
Suppose, for contradiction, that membership in \(S_p\) were ultimately periodic with period \(h\geq1\). Then \(Mh\) would also be an ultimate period. For all sufficiently large \(q\), equation~\eqref{eq:rail-p} would give
\[
q\in Y_0
\iff Mq+p\in S_p
\iff M(q+h)+p\in S_p
\iff q+h\in Y_0,
\]
contradicting Lemma~\ref{lem:Y0-aperiodic}. Thus the membership indicator of \(S_p\) is not ultimately periodic. The contrapositive of Lemma~\ref{lem:gaps-imply-membership} now shows that the consecutive gaps in the increasing enumeration of \(S_p\) are not eventually periodic. By Lemma~\ref{lem:greedy}, this enumeration is the greedy extension of \(A_p\), proving the theorem.
\end{proof}

\begin{thebibliography}{9}

\bibitem{ErdosGraham1980}
P. Erd\H{o}s and R.~L. Graham, \emph{Old and New Problems and Results in Combinatorial Number Theory}, Monographies de L'Enseignement Math\'ematique, no.~28, Universit\'e de Gen\`eve, L'Enseignement Math\'ematique, Geneva, 1980.

\bibitem{Bloom341}
T.~F. Bloom, \emph{Erd\H{o}s Problem \#341}, Erd\H{o}s Problems, \href{https://www.erdosproblems.com/341}{\texttt{erdosproblems.com/341}}, accessed 11 August 2026.

\bibitem{CalkinFinch1996}
N.~J. Calkin and S.~R. Finch, \emph{Conditions on periodicity for sum-free sets}, Experiment.\ Math.\ \textbf{5} (1996), no.~2, 131--137; \href{https://doi.org/10.1080/10586458.1996.10504584}{doi:10.1080/10586458.1996.10504584}.

\bibitem{CalkinErdos1996}
N.~J. Calkin and P. Erd\H{o}s, \emph{On a class of aperiodic sum-free sets}, Math.\ Proc.\ Cambridge Philos.\ Soc.\ \textbf{120} (1996), no.~1, 1--5; \href{https://doi.org/10.1017/S0305004100074600}{doi:10.1017/S0305004100074600}.

\end{thebibliography}

\end{document}
